\section{怎样使用这一章：先把方程真正算出来}
\label{sec:tutorial-start}
若希望先了解核心方程、模型变种、实现流程及实验适用性，可先跳到第\ref{ch:core-guide}章；本章则负责把每一步推导和模型假设展开。
本章按可以拿纸笔跟算的讲义组织。学习目标不是背下几个模型名，而是能够回答：未知量有哪些，方程从哪里来，哪一步使用了近似，删掉某一项以后还允许什么物理过程。前半章建立共同基础，后半章用同一套工具解释具体论文。遇到专题里暂时不熟悉的参数，可以返回本章的定义和量纲表；不必第一次阅读就掌握全部专题。

\subsection{建议分六遍学习，每遍留下一个能独立完成的结果}
\begin{longtable}{P{14mm}P{66mm}P{69mm}}
\toprule 次序 & 阅读任务 & 合上讲义后应能完成的事\\\midrule
1 & 分布函数、速度积分、连续性和动量矩 & 不看答案推出连续性和动量方程，指出分部积分的边界条件\\
2 & 压力张量、内能、五矩和十矩 & 展开三阶原始矩；解释二维空间为何仍可有六个压力分量\\
3 & Maxwell约束、质心变换和欧姆定律 & 正确写出电子负电荷的符号；区分离子速度与质心速度\\
4 & 冻结、能量、Hall及惯性尺度 & 证明Hall项总功为零；说明为什么这不等于没有能量转换\\
5 & EMHD、强导引场和两场模型 & 手算一个圆偏振色散关系，验证两场能量的耦合项抵消\\
6 & 撕裂、FLR、闭合及文献专题 & 明确外区、内区和匹配条件；逐篇列出模型额外假设\\\bottomrule
\end{longtable}
每次宜先跟算一个完整小节，再做附近的验算。后面的专题不要求在首次学习时同时读完。所有推导仍集中在第一章；算法实现与文献结果放在后续章节。

\subsection{坐标、下标与单位：先消除最容易混淆的地方}
本章基础部分取$x$为流出、$y$为流入、$z$为电流和导引场方向。二维表示$\partial_z=0$，并不表示$u_z,B_z,E_z$为零。物种下标$s=e,i$分别表示电子和离子；张量中的$i,j,k$只是笛卡尔分量指标。若两种$i$同时出现，应先按上下文区分物种和分量。

默认单价、完全电离、非相对论电子--离子等离子体，$q_i=e>0$，$q_e=-e$。温度$T_s$用能量单位，例如$1\,{\rm eV}=1.602176634\times10^{-19}\,{\rm J}$；因此$p_s=n_sT_s$不再额外乘$k_B$。无量纲方程会明确说明基准量。在部分论文专题中需要更换坐标或用$\rho_s$表示声回旋半径，会在该处重新定义；基础部分的$\rho_s=m_sn_s$只表示质量密度，声回旋半径写$\rho_{\rm sound}$。

\begin{longtable}{P{31mm}P{61mm}P{57mm}}
\toprule 记号 & 定义或意义 & SI单位\\\midrule\endfirsthead
\toprule 记号 & 定义或意义 & SI单位\\\midrule\endhead
$n_s,\rho_s$ & 数密度，质量密度$m_sn_s$ & $ {\rm m^{-3}},\,{\rm kg\,m^{-3}}$\\
$\vel_s,\vct c$ & 平均速度，随机速度$\vct v-\vel_s$ & $ {\rm m\,s^{-1}}$\\
$P_{s,ij}$ & 随机运动的动量通量 & $ {\rm Pa=J\,m^{-3}}$\\
$Q_{s,ijk},q_{s,k}$ & 三阶中心矩，热能通量向量 & $ {\rm W\,m^{-2}}$\\
$\eta,D_\eta$ & 电阻率，磁扩散率$\eta/\mu_0$ & $\Omega\,{\rm m},\,{\rm m^2\,s^{-1}}$\\
$d_i,d_e$ & 离子、电子惯性长度 & $ {\rm m}$\\
$\omega,\gamma$ & 波频率、模增长率 & $ {\rm s^{-1}}$\\\bottomrule
\end{longtable}
特别注意：电阻率与磁扩散率的单位不同；热流张量和热流向量虽然单位相同，却不是同一个量。温度张量常被定义为$\mathsf T_s=\mathsf P_s/n_s$，也有代码再除以质量，比较公式时必须查清定义。

\subsection{张量计算的最低准备}
重复分量指标求和，例如$a_ib_i=a_xb_x+a_yb_y+a_zb_z$。并矢和双重缩并分别表示
\begin{equation}
 (\vct a\vct b)_{ij}=a_ib_j,\qquad
 \mathsf P:\grad\vel=\sum_{i,j}P_{ij}\partial_ju_i.
\end{equation}
压力散度的$x$分量是
\begin{equation}
 (\grad\cdot\mathsf P)_x=\partial_xP_{xx}+\partial_yP_{xy}+\partial_zP_{xz}.
\end{equation}
它一般不是$\partial_xp$。只有先假定$P_{ij}=p\delta_{ij}$，非对角项才为零，三个对角项才相等。压力为正也不只是要求三个对角元为正：应有$\vct a\cdot\mathsf P\cdot\vct a\ge0$对任意$\vct a$成立。

随体导数$D_s=\partial_t+\vel_s\cdot\grad$表示沿该物种平均流轨迹看到的变化。对任意标量$a$，乘积法则给出
\begin{align}
 \partial_t(na)+\grad\cdot(n\vel a)
 &=n(\partial_ta+\vel\cdot\grad a)
   +a[\partial_tn+\grad\cdot(n\vel)]\\
 &=nDa,\qquad\text{最后一步才使用连续性。}
 \label{eq:tutorial-product}
\end{align}
后面所有“守恒形式转为随体形式”，都可以用这个恒等式逐分量完成。它不要求不可压；不可压$\grad\cdot\vel=0$与密度随流保持$Dn=0$的关系，还必须通过连续性来判断。

\subsection{模型地图：每条箭头都要带着条件}
共同起点是粒子分布的输运。取矩得到无穷方程链，选闭合才得到有限流体；准中性、忽略位移电流、忽略电子惯性，又是另外三类近似。它们不能合并为一句“取低频极限”。
\begin{longtable}{P{30mm}P{59mm}P{60mm}}
\toprule 模型 & 主要未知量 & 进入该模型时额外规定什么\\\midrule
完整双流体 & 两物种密度、流速、压力矩及$\E,\B$ & 各物种的热流或更高矩闭合\\
五矩 / 十矩 & 每物种5或10个矩 & 标量或张量压力及相应闭合\\
准中性双流体 & 共同密度、两流体速度和压力 & 长尺度、低频及准中性排序\\
Hall MHD & 总密度、离子动量、磁场和压强 & 忽略电子惯性，保留相对漂移\\
惯性扩展MHD & 质心量、电流或广义磁场 & 惯性、漂移应力和能量一致\\
EMHD & 给定离子背景，演化电子流和磁场 & 忽略所研究尺度的离子响应\\
约化 / 陀螺流体 & 磁通、流函数及选定矩 & 强磁化低频排序及FLR近似\\\bottomrule
\end{longtable}
例如五矩含有完整电磁波，却没有十矩的非对角压力；十矩含有压力各向异性，却没有完整粒子分布。模型不是按名称排成一条从“差”到“好”的直线，应由目标尺度和需要检验的现象来选择。\citep{Hakim2008,Wang2015}
